
\medskip
\noindent{\it Exercise \the\chapternum.13}\quad A household has preferences over
consumption processes $\{c_t\}_{t=0}^\infty$ that are  ordered by
$$ -  .5 \sum_{t=0}^\infty \beta^t [(c_t - 30)^2  + .000001b_t^2 ] \eqno(1)$$
where $\beta=.95$.  The household chooses a consumption, borrowing
plan to maximize (1) subject to the sequence of budget constraints
$$ c_t + b_t = \beta b_{t+1} + y_t  $$ % \EQN e132 $$
for $t \geq 0$, where $b_0$ is an initial condition,
$\beta^{-1}$ is the one-period gross risk-free interest rate,
$b_t$ is the household's one-period debt that is due in period $t$,
and $y_t$ is its labor income, which obeys the second-order autoregressive
process
$$ (1-\rho_1 L - \rho_2 L^2)y _{t+1}  =    (1 -\rho_1 -\rho_2) 5 +
    .05 w_{t+1} $$ % \EQN 133 $$
where $\rho_1 =1.3, \rho_2 = -.4$.

\medskip
\noindent{\bf a.}  Define the {\it state\/} of the household  at $t$ as
$x_t= \left[\matrix{1 & b_t & y_t & y_{t-1} \cr}\right]'$ and
the {\it control\/} as   $u_t = (c_t -30)$.  Then express the
transition law facing the household in the form \Ep{lasset1}.
Compute the eigenvalues of $A$.  Compute the
zeros of the characteristic polynomial $(1-\rho_1 z - \rho_2 z^2)$ and  compare
them with the eigenvalues of $A$. ({\it Hint:} To compute the zeros in
 Matlab,
set
%$a=\left[ \matrix{1 &  -1.3  & .4\cr}\right]$
$a=\left[ \matrix{.4 &  -1.3  &  1 \cr}\right]$
 and call
{\tt roots(a)}. The zeros of $(1 -\rho_1 z  - \rho_2 z^2) $ equal
the {\it reciprocals\/} of the eigenvalues of the associated $A$.)

\medskip
\noindent{\bf b.}  Write a Matlab program that uses the Howard improvement
algorithm \Ep{howardimprove} to compute the household's optimal decision rule
for $u_t = c_t -30$.   Tell how many iterations it takes for this to converge
(also tell your convergence criterion).
 \medskip
\noindent{\bf c.}  Use the household's optimal decision rule
to compute the law of motion for $x_t$ under the optimal decision rule in
the form $$ x_{t+1} = (A - BF^*) x_t + C w_{t+1}, $$ where $u_t = - F^* x_t$
is the optimal decision rule.  Using Matlab,
compute the impulse response function
of $\left[\matrix{c_t & b_t \cr}\right]'$ to $w_{t+1}$.
Compare these with the theoretical expressions
 \Ep{sprob16}.
 
 
\medskip

\noindent{\it Exercise \the\chapternum.19}\quad  {\bf IQ}
\medskip
\noindent An infinitely lived  person's `true intelligence'  $\theta \sim {\cal N}(100, 100)$, i.e., mean 100, variance 100. For each date $t \geq 0$,
the person takes a `test' with the  outcome being a univariate random variable   $y_t = \theta + v_t$,
where $v_t$ is an i.i.d.~process with distribution ${\cal N}(0, 100)$. The person's initial IQ is $IQ_0=100$ and at date  $t\geq 1$ before the date $t$ test is
taken it is ${\rm IQ}_t = E \theta | y^{t-1}$, where $y^{t-1}$ is the history of test scores from date $0$ until date $t-1$.

\medskip
\noindent{\bf a.}  Give a recursive formula for ${\rm IQ}_t$ and for $E ({\rm IQ}_t - \theta)^2$.
\medskip
\noindent {\bf b.}  Use Matlab to simulate 10 draws of $\theta$ and associated paths of $y_t, {\rm IQ}_t$ for $t=0, \ldots, 50$.
\medskip
\noindent {\bf c.}  Prove that $  \lim_{t\rightarrow \infty} E ({\rm IQ}_t - \theta)^2=0$.

\medskip

\noindent{\it Exercise \the\chapternum.20}\quad  {\bf Random walk}
\medskip
\noindent A scalar process $x_t$ follows the process
$$ x_{t+1} = x_t + w_{t+1}  $$
where $w$ is an i.i.d.~${\cal N}(0,1)$ scalar process and $x_0 \sim {\cal N}(\hat x_0, \Sigma_0)$.
Each period, an observer receives two signals in the form of a $2 \times 1$ vector $y_t$ that obeys
$$ y_t = \bmatrix{1 \cr 1\cr} x_t + v_t $$
where the $2 \times 1$ process $v_t$ is i.i.d.~with distribution $v_t \sim {\cal N}(0, R)$ where
$R =\bmatrix{ 1 & 0 \cr 0 &1  \cr}$.

\medskip
\noindent{\bf a.}  Suppose that $\Sigma_0 = 1.36602540378444$.  For $t \geq 0$, find formulas for
$E [x_t |  y^{t-1}] $, where $y^{t-1}$ is the history of $y_s$ for $s$ from $0$ to $t-1$.

\medskip
\noindent{\bf b.} Verify numerically that the matrix $A-KG$ in formula \Ep{VAR1} is stable.

\medskip
\noindent{\bf c.}  Find an infinite-order vector autoregression for $y_t$.




\medskip

\noindent{\it Exercise \the\chapternum.23}\quad  {\bf A monopolist, learning, and ergodicity}
\medskip
 \noindent A monopolist  produces a quantity $Q_t$ of a single good in every period $t \geq 0$ at zero cost.
 At the beginning of each period $t \geq 0$, before output price $p_t$ is observed, the monopolist sets quantity $Q_t$ to maximize
$$ E_{t-1} p_t Q_t \leqno(1)  $$
where $p_t$ satisfies the linear inverse demand curve
$$ p_t = a - b Q_t + \sigma_p \epsilon_{t}  \leqno(2) $$
where $b>0$ is a constant known to the firm,  $\epsilon_t$ is an i.i.d.~scalar with distribution $\epsilon_t \sim {\cal N}(0,1)$, and the constant
in the inverse demand curve
$a$ is a  scalar random variable unknown to the firm and whose  unconditional distribution is $a \sim {\cal N}(\mu_a, \sigma_a^2)$,
where $\mu_a>0$ is large relative to $\sigma_a>0$. Assume that the random variable $a$  is independent of $\epsilon_t$ for all $t$. Before the firm chooses
$Q_0$, it knows the unconditional distribution of $a$, but not the realized value of $a$. For each $t \geq 0$, the firm wants to estimate  $a$ because it wants to make a good decision about output $Q_t$.
 At the end of each period $t$, when it must set $Q_{t+1}$, the firm observes $p_t$ and also of course knows the value of $Q_t$ that it had
 set.
 In (1), for $t \geq 1$,  $E_{t-1} (\cdot) $ denotes the mathematical expectation conditional on the history of signals $p_s, q_s, s= 0, \ldots, t-1$;
for $t=0$, $E_{t-1} (\cdot)$ denotes the expectation conditioned on no previous  observations of $p_t, Q_t$.

\medskip
\noindent {\bf a.}    What is the optimal setting for $Q_0$?  For each date $t \geq 0$, determine the firm's optimal setting for  $Q_t$ as a function
of the information $p^{t-1}, Q^{t-1}$ that the firm has when it sets $Q_t$.
\medskip
\noindent{\bf b.}  Under the firm's optimal policy, is  the pair $(p_t, Q_t)$ Markov?
\medskip
\noindent{\bf c.} `Finding the state is an art.' Find a recursive representation of the firm's optimal policy for setting $Q_t$ for $t \geq 0$.  Interpret the state variables that you
propose.
\medskip
\noindent  {\bf d.} Under the firm's optimal rule for setting $Q_t$, does  the random variable $ E_{t-1} p_t$  converge to a constant
as $t \rightarrow +\infty$?  If so, prove that it does and find the limiting value. If not, tell why it does not converge.
\medskip
\noindent {\bf e.}  Now suppose that instead of maximizing (1) each period, there is a single infinitely lived monopolist who
once and for all before time $0$ chooses a plan for an entire sequence  $\{Q_t\}_{t=0}^\infty$, where the $Q_t$ component has to be a measurable function
of $(p^{t-1}, q^{t-1})$, and where the monopolist's objective is to maximize
$$ E_{-1} \sum_{t=0}^\infty \beta^t p_t Q_t \leqno(3)  $$
where $\beta \in (0,1)$ and $E_{-1}$ denotes the mathematical expectation conditioned on the null history. Get as far as you can
in deriving the monopolist's optimal sequence of decision rules.


\medskip
\noindent{\it Exercise \the\chapternum.25}\quad  {\bf Consumption}
\medskip
\noindent {\bf a.}  Please use formulas \Ep{sprob16} to verify formulas  \Ep{consexample1}
and \Ep{consexample2}-\Ep{incomemaar} of subsection \use{sec:classic_consumption}.
\medskip
\noindent{\bf b.} Please use formulas \Ep{muth2} to compute the decision rules
in formulas   \Ep{consexample1} and \Ep{consexample2} for the following parameter values:
$\beta = .95, \sigma_1= \sigma_2 =1$.
\medskip
\noindent{\bf c.}
 Please use formulas \Ep{muth2} to compute the decision rules
in formulas   \Ep{consexample1} and \Ep{consexample2} for the following parameter values:
$\beta = .95, \sigma_1= 2,  \sigma_2 =1$.

\medskip
\noindent{\bf d.} Please use formula \Ep{debt_evolution} to confirm formulas
\Ep{consexample1a} and \Ep{consexample1b}.

\medskip

\noindent{\it Exercise \the\chapternum.26}\quad  {\bf Math and verbal IQ's}
\medskip
\noindent An infinitely lived  person's `true intelligence' $\theta$ has two components, math ability $\theta_1$ and
verbal ability $\theta_2$, where  $\theta \sim {\cal N}\Biggl(\bmatrix{100 \cr 100},  \bmatrix{100 & 0 \cr 0 & 100}\Biggr)$.
 For each date $t \geq 0$,
the person takes a single `test' with the  outcome being a univariate random variable   $y_t = G_t \theta + v_t$,
where $v_t$ is an i.i.d.~process with distribution ${\cal N}(0, 50)$ and $G_t =\bmatrix{.9 & .1 }$ for $t = 0, 2, 4, \ldots$ and
$G_t = \bmatrix{.01 & .99}$ for $t=1, 3, 5, \ldots$.  Here the person takes a math test at $t$ even and a verbal test at $t$ odd (but you have to know
how to read English to survive the math test, and you have to know how to tell time in order to plan your time allocation well for the verbal test).
 The person's initial IQ vector is $IQ_0=\bmatrix{100 \cr 100}$ and at date  $t\geq 1$ before the date $t$ test is
taken it is ${\rm IQ}_t = E \theta | y^{t-1}$, where $y^{t-1}$ is the history of test scores from date $0$ until date $t-1$.

\medskip
\noindent{\bf a.}  Give a recursive formula for ${\rm IQ}_t$ and for $E ({\rm IQ}_t - \theta)({\rm IQ}_t - \theta)'$.
\medskip
\noindent {\bf b.}  Use Matlab to simulate 10 draws of $\theta$ and associated paths of $y_t, {\rm IQ}_t$ for $t=0, \ldots, 50$.
\medskip
\noindent {\bf c.} Show computationally or analytically  that $\lim_{t\rightarrow +\infty}E({\rm IQ}_t - \theta)({\rm IQ}_t - \theta)'=\bmatrix{0 & 0 \cr 0 & 0}$.




\medskip

\noindent{\it Exercise \the\chapternum.27}\quad  {\bf Permanent income model again}
\medskip
\noindent  Each of two consumers named $i=1,2$  has preferences over consumption streams
that are ordered by
the utility functional
$$  E_0 \sum_{t=0}^\infty \beta^t u(c_t^i)  \leqno(1) $$
where $E_t$ is the mathematical expectation conditioned
on the consumer's time $t$ information,  $c_t^i$ is time $t$ consumption of consumer $i$ at time $t$,
$u(c)$ is a strictly concave one-period utility function, and
$\beta \in (0,1)$ is a discount factor.  The consumer maximizes
(1)  by choosing a consumption, borrowing plan
 $\{c_t^i, b_{t+1}^i\}_{t=0}^\infty$ subject to the sequence of budget constraints
$$ c_t^i + b_t^i = R^{-1} b_{t+1}^i  + y_t^i  $$
where $y_t$ is an exogenous
 stationary endowment process, $R$ is a constant gross
risk-free interest rate, $b_t^i$ is one-period risk-free  debt  maturing at
$t$, and $b_0^i = 0$ is a given initial condition.  Assume
that $R^{-1} = \beta$. We impose the following condition on the
consumption, borrowing plan of consumer $i$:
$$  % \lim_{T \rightarrow +\infty}
 E_0 \sum_{t=0}^\infty \beta^t (b_t^i)^2 < +\infty. $$ % \EQN sprob3 $$
Assume
the quadratic  utility function
$u(c_t) =  -.5 (c_t - \gamma)^2$,
where $\gamma > 0 $ is a bliss level of consumption. Negative consumption  rates
are allowed.

 Let $s_t \in \{0,1\}$ be an i.i.d.~process
with ${\rm Prob} (s_t =1) = {\rm Prob} (s_1 = 0 ) = .5$.  The endowment process
of consumer 1 is $y_t^1 = 1 - .5 s_t$ and the endowment process of person 2 is
$y_t^2 = .5+ .5 s_t$.  Thus, the two consumers' endowment processes are perfectly
negatively correlated i.i.d.~processes with means of .75.


\medskip
\noindent{\bf a.} Find  optimal decision rules for consumption for both consumers.
Prove that the consumers' optimal decisions imply the following laws of motion
for $b_t^1, b_t^2$:
$$ \eqalign{ b_{t+1}^1(s_t =0) & = b_t^1 - .25 \cr
            b_{t+1}^1(s_t =1) & = b_t^1 + .25 \cr
            b_{t+1}^2(s_t =0) & = b_t^2 + .25 \cr
            b_{t+1}^2(s_t =1) & = b_t^2 - .25 \cr } $$
\medskip
\noindent{\bf b.} Show that for each consumer, $c_t^i, b_t^i$ are co-integrated.
\medskip
\noindent{\bf c.} Verify that $b_{t+1}^i$ is risk-free in the sense that conditional on
information available at time $t$, it is independent of news arriving at time $t+1$.
\medskip
\noindent{\bf d.} Verify that with the initial conditions $b_{0}^1 = b_0^2 =0$, the following
two equalities obtain:
$$ \eqalign{ b_t^1 + b_t^2& = 0 \quad \forall t \geq 1 \cr
             c_t^1 + c_t^2 & = 1.5 \quad \forall t \geq 1 \cr} $$
Use these conditions to interpret the decision rules that you have computed as describing a closed pure consumption loans economy
in which consumers $1$ and $2$ borrow and lend with each other and in which the risk-free asset is a one-period IOU from one of the consumers to the other.
\medskip

\noindent{\bf e.}  Define the `stochastic discount factor of consumer $i$' as
$m_{t+1}^i = {\frac{\beta u'(c_{t+1}^i)}{u'(c_t^i)}}$.  Show that the stochastic discount factors of consumer $1$ and $2$ are
$$  m_{t+1}^1 = \cases{ \beta + .25 {\frac{\beta(1-\beta)}{(\gamma- c_t^1)}} , & if $s_{t+1} = 0 $; \cr
                          \beta - .25 {\frac{\beta(1-\beta)}{(\gamma- c_t^1)}} , & if $s_{t+1} = 1 $;    . \cr}  $$
$$  m_{t+1}^2 = \cases{ \beta - .25 {\frac{\beta(1-\beta)}{(\gamma- c_t^2)}} , & if $s_{t+1} = 0 $; \cr
                          \beta + .25 {\frac{\beta(1-\beta)}{(\gamma- c_t^2)}} , & if $s_{t+1} = 1 $;    . \cr}  $$
Are the stochastic discount factors of the two consumers equal?
\medskip
\noindent{\bf f.} Verify that $E_t m_{t+1}^1 = E_t m_{t+1}^2 = \beta$.
\medskip



\vfill\break


\medskip

\noindent{\it Exercise \the\chapternum.28}\quad  {\bf Invertibility}
\medskip
\noindent A univariate stochastic process $y_t$ has a first-order moving
average representation
$$ y_t = \epsilon_t - 2 \epsilon_{t-1} \leqno(1) $$
where $\{\epsilon_t\}$ is an i.i.d.~process distributed
${\cal N}(0,1)$.

\medskip
\noindent{\bf a.}  Argue that $\epsilon_t$ cannot be expressed as as linear combination of
$y_{t-j}, j \geq 0$ where the sum of the squares of the weights is finite.  This means
that $\epsilon_t$ is not in the space spanned by square summable linear combinations
of the infinite history $y^t$.
\medskip
\noindent{\bf b.}  Write equation (1) as a state space system, indicating the matrices
$A, C, G$.
\medskip
\noindent {\bf c.} Using the matlab program {\tt kfilter.m} \mtlb{kfilter.m}%
to compute an innovations representation for $\{y_t\}$.  Verify that the innovations
representation for $y_t$ can be represented as
$$ y_t = a_t - .5 a_{t-1} \leqno(2) $$
where $a_t = y_t - E [y_t| y^{t-1} ] $ is a serially uncorrelated process.
Compute the variance of $a_t$.  Is it larger or smaller than the variance of $\epsilon_t$?

\medskip
\noindent{ \bf d.}  Find an autoregressive representation for $y_t$ of the form
$$ y_t = \sum_{j=1}^\infty A_j y_{t-j} + a_t \leqno(3) $$
where $E a_t y_{t-j} = 0 $ for $j \geq 1$. ({\it Hint:} either use formula (2) or else  remember formula \Ep{VAR1}.)

\medskip
\noindent{\bf e.}  Is $y_t$ Markov?  Is $\bmatrix{y_t & y_{t-1}}'$ Markov? Is $\bmatrix{y_t & y_{t-1} & \cdots & y_{t-10}}'$ Markov?

\medskip
\noindent{\bf f.}  Extra credit.  Verify that $\epsilon_t$ {\it can\/} be expressed as a square summable linear combination
of $y_{t+j}, j\geq 1$.


%
\medskip
\noindent {\it Exercise \the\chapternum.33}\quad  {\bf From shocks to innovations, I}

\medskip
\noindent
Consider again the setting described in exercise  {\it \the\chapternum.32\/}.  Let $\{y_t\}_{t=-\infty}^\infty$ be the covariance stationary scalar
stochastic process appearing in exercise {\it \the\chapternum.32\/}.  We begin by introducing a
covariance generating function defined as the $z$-transform of the autocovariance sequence $\{c_y(\tau)\}_{\tau =-\infty}^\infty$,
where $c_y(\tau) = E y_t y_{t-\tau}$.  The covariance generating function is
$$ S_y(z) = \sum_{j=\infty}^\infty c_y(j) z^j.$$  Using the original state-space representation for $y_t$, it can be verified that
$$ \eqalign{ S_y(z) & = r + \left({\frac{c}{1-\rho z}}\right) \left({\frac{c}{1-\rho z^{-1}}} \right) \cr
                    & = {\frac{c^2 + r(1-\rho z)(1 - \rho z^{-1} ) }{(1 -\rho z)(1 -\rho z)}} \cr
                    & = {\frac{ c^2 + r(1 + \rho^2) - r \rho z - r \rho z^{-1}}{(1 -\rho z)(1 -\rho z^{-1})}}. } $$
(The spectral density defined in section \use{sec:spectrum} is the covariance generating function obtained by setting $z= e^{-i \omega}, \omega \in [-\pi, \pi]$.)
 Using  the innovations representation for $\{y_t\}$ instead of the original state-space representation cast in terms of the hidden state $x_t$,
 it can  be verified that the covariance generating function also equals
$$ S_y(z) =  {\frac{\omega ( 1 - (\rho -k )z)( 1 - (\rho - k)z^{-1} )}    {(1 -\rho z)(1 -\rho z^{-1})}} .  $$
The denominators in both representations of $S_y(z)$ are equal and the numerators in both representations must also be equal.  Therefore,  it must be  true that
$$   c^2 + r(1 + \rho^2) - r \rho z - r \rho z^{-1} = \omega ( 1 - (\rho -k )z)( 1 - (\rho - k)z^{-1}) . \leqno(1)  $$
The left side is a  quadratic  polynomial in $z$. The right side factors this polynomial. The
zeros of the quadratic polynomial  in $z$ evidently form  a reciprocal pair  $(\rho -k), (\rho -k)^{-1}$.\NFootnote{Such a reciprocal pairs property
characterizes large classes of linear filtering problems and linear optimal control problems.  We shall encounter more examples in chapter
\use{dplinear}.}  Equation   (1) as an instance of a more general   {\it  spectral factorization identity\/} (see Hansen and Sargent (2013, p.~164).)
 Also, please note the steady state Kalman filter generates a Kalman
gain $k$ that is a key to  achieving the factorization.
\index{spectral factorization identity}%
\auth{Hansen, Lars Peter}%
\auth{Sargent, Thomas J.}%

Now take representation ($\dagger$) for $a_t$ from exercise {\it \the\chapternum.32\/} and use it to compute the covariance generating
function of the innovation process $\{a_t\}$:
$$ \eqalign{ S_a(z) & = {\frac{c^2}{(1 - (\rho - k) z) (1 - (\rho - k) z^{-1})}} + {\frac{r (1 - \rho z) (1 - \rho z^{-1})}{(1 - (\rho -k)z)
         ( 1 - (\rho - k) z^{-1}) }} \cr
         & = {\frac{  c^2 + r(1 + \rho^2) - r \rho z - r \rho z^{-1}  }  { ( 1 - (\rho -k )z)( 1 - (\rho - k)z^{-1})  }} .  } $$
It follows from the preceding equation and  factorization identity (1) that
$$ S_a(z) = \omega,  $$
which implies that $c_a(j) = 0 $ for $j \neq 0$ and $c_a(0) = \omega$. This confirms that   $\{a_t\}$ is a serially uncorrelated process with variance $\omega$.

\medskip
\noindent {\bf a.} Please verify the preceding calculations.

\medskip
\noindent{\bf b.}  Please plot the spectral density $S_a(e^{-i \omega})$ against $\omega \in [0, \pi]$ and compare it with those of the scalar processes
depicted in   Figures \Fg{r2red1a}, \Fg{r2red1b}, \Fg{r2red1d}, and \Fg{r2red1c} from section \use{sec:spectrum}.






\medskip
\noindent {\it Exercise \the\chapternum.34}\quad  {\bf From shocks to innovations, II}


\medskip
\noindent Consider a state space system and the associated time-invariant innovations representation,
which we repeat here for convenience. The original system is
$$ \eqalign{ x_{t+1} & = A x_t + C w_{t+1} \cr
         y_t & = G x_t  + v_t } \leqno(1)  $$
where $\{w_{t+1}\}_{t=0}^\infty $ is an i.i.d.~sequence with $w_{t+1} \sim {\cal N}(0,I)$ for all $t$
and $\{v_t\}_{t=0}^\infty$ is an i.i.d.~sequence with $v_t \sim {\cal N}(0, R)$ for all $t$; and $E w_t' v_s = 0$
for all $t$ and $s$.
The associated time-invariant ``innovations representation'' state-space system is
$$ \eqalign{ \hat x_{t+1} & = A \hat x_t + K a_t \cr
              y_t & = G \hat x_t + a_t. } \leqno(2) $$
Here $C w_{t+1}$ are the shocks hitting the hidden state $x_{t+1}$ conditional on  $x_t$,
$v_t$ are shocks or measurements hitting $y_t$ conditional on $G x_t$, and $a_t = y_t - E [y_t | y^{t-1}]  $ are
shocks hitting $y_t$ given the infinite history $y^{t-1} = [ y_{t-1}, y_{t-2}, \ldots ]$.
The shock $a_t$ is i.i.d.~with distribution $a_t \sim {\cal N} (0, \Omega)$.

\medskip
\noindent{\bf a.} Please provide formulas for $K$ and $\Omega$.

\medskip

\noindent{\bf b.} Please deduce from the above pair  of state-space systems a single state-space  system that maps the
shocks $w_{t+1}, v_{t+1}$ into the innovation $a_{t+1}$ and verify that it takes the form:

$$ \eqalign{ \bmatrix{x_{t+1} \cr \hat x_t \cr v_{t+1} } & = \bmatrix{A & 0 & 0 \cr
                                                                    K G & (A-KG) &  K \cr
                                                                    0 & 0 & 0 }
                                                           \bmatrix{x_t \cr \hat x_t \cr v_t } +
                                                           \bmatrix{ C & 0 \cr
                                                             0 & 0 \cr
                                                             0 & I }
                                                             \bmatrix{ w_{t+1} \cr v_{t+1} }  \cr
              a_t & = \bmatrix{ G & - G & I } \bmatrix{ x_t \cr \hat x_t \cr v_t } }    \leqno(3)                                             $$

\medskip

\noindent {\bf c.}  Please use formulas from section \use{sec:stochlinear}    to compute moving average representations and autocovariograms
for the process $\{a_t\}$.
Is $\{a_{t}\}$ a serially correlated process?


\medskip
\noindent{\bf d.}  Please use  formulas from part {\bf c.} to deduce impulse response functions of $a_t$ to $w_t$ and of
$a_t$ to $v_t$ so that $a_t = \sum_{j=0}^\infty g_j w_{t-j} +  \sum_{j=0}^\infty h_j v_{t-j}$.  Please represent the
autocovariance function of $\{a_t\}$ as the sum of the autocovariance functions of $\sum_{j=0}^\infty g_j w_{t-j} $
and $\sum_{j=0}^\infty h_j v_{t-j}$ and  verify that $E a_t a_{t-j}' = 0 $ for $j\neq 0$. (Feel free to write a Matlab or Python program
to perform these calculations and verifications.)




\medskip
\noindent {\it Exercise \the\chapternum.36}\quad  {\bf IQ again}

\medskip
\noindent Latent intelligence is a scalar random variable $\theta \sim {\cal N}(\mu_\theta, \sigma_\theta^2)$.  An applicant for college
takes $k$ tests whose scores are a $k\times 1$ random vector governed by
$$ y = \lambda \theta + u \leqno(1) $$
where $\lambda$ is a known $k\times 1 $ vector,  $u \sim {\cal N}(0, \Sigma_u)$, and $E u \theta =0$.
% Sometimes the covariance matrix $\Sigma_u$ is
diagonal, but that is not necessary.
  The mean and variance parameters $\mu_\theta, \sigma_\theta$ and the covariance matrix
$\Sigma_u$ as well as the vector of ``loadings'' $\lambda$ are all known.  A testing service summarizes results of the
test as a single score $IQ = E [\theta | y]$, the mathematical expectation of the student's  $\theta$ conditional on his/her test scores $y$.
Please show that $IQ$ is an affine  function of $y$ and give formulas for the constant and slope coefficients as functions
of the known parameters $\lambda, \mu_\theta, \sigma_\theta, \Sigma_u$.

\vfil\eject







